\providecommand{\Inn}{\operatorname{Inn}}
\providecommand{\PSU}{\operatorname{PSU}}
\providecommand{\Sz}{\operatorname{Sz}}
\providecommand{\Ree}{\operatorname{Ree}}
\providecommand{\Suz}{\operatorname{Suz}}
\providecommand{\G}{\operatorname{G}}

\chapter{Michael Aschbacher (2012)}
\index{Aschbacher, Michael}

\begin{quote}
\it ``The classification of the finite simple groups stands as one of the greatest achievements of twentieth-century mathematics. Michael Aschbacher has been the dominant figure in the later stages of this project, and his work has been essential for bringing it to completion.'' --- Wolf Prize Citation, 2012
\end{quote}

\section{Main Contribution}

Michael Aschbacher (born 1944) is an American mathematician renowned for his monumental contributions to the classification of finite simple groups. He was awarded the Wolf Prize in Mathematics in 2012 alongside Luis A. Caffarelli.

The Wolf Prize citation reads: ``for his work in classifying finite simple groups.''

Aschbacher's core contributions include:
\begin{itemize}
    \item \textbf{The Classification of Finite Simple Groups:} A central role in completing the largest and most complicated portion of the classification, including the analysis of quasithin groups, groups of characteristic 2-type, and the revisionist program.
    \item \textbf{Aschbacher's Theorem:} A classification of the maximal subgroups of finite classical groups, which is a fundamental tool in group theory and its applications.
    \item \textbf{Aschbacher--Guralnick Theorem:} A theorem on the generation of finite simple groups, with important consequences for the structure of group extensions.
    \item \textbf{Contribution to the Monster Group:} Determining the existence and uniqueness of the largest sporadic simple group (the Monster), including the identification of the Fischer--Griess Monster.
    \item \textbf{Theory of Fusion Systems:} Pioneering work on fusion systems, which has become a major area of research in its own right.
\end{itemize}

\section{Origin of the Problem}

\subsection{The Classification Problem}

A finite simple group is a nontrivial finite group with no proper nontrivial normal subgroups. By the Jordan--H\"older theorem, every finite group can be built from simple groups in an essentially unique way. Thus classifying all finite simple groups became the central problem of finite group theory.

By the early 1970s, several infinite families of finite simple groups were known:
\begin{itemize}
    \item The cyclic groups of prime order, $\Z/p\Z$.
    \item The alternating groups $A_n$ for $n \geq 5$.
    \item The groups of Lie type --- Chevalley groups, twisted groups (Steinberg, Suzuki, Ree).
\end{itemize}
In addition, a number of sporadic simple groups had been discovered (the Mathieu groups, Janko groups, etc.).

The classification project aimed to prove that every finite simple group belongs to one of these families or is one of the finitely many sporadic groups.

\subsection{The Size and Complexity}

The CFSG is estimated to span thousands of journal pages spread over hundreds of papers by dozens of mathematicians. By the 1980s, Aschbacher emerged as the key figure managing the most difficult parts, particularly:
\begin{itemize}
    \item The classification of quasithin groups (groups of low 2-rank).
    \item The analysis of groups of characteristic 2-type.
    \item The revision and consolidation of earlier partial results.
\end{itemize}

\subsection{The Revisionist Program}

As Aschbacher recognized, the original proof of the classification was sprawling, with overlaps, gaps, and errors. He initiated the revisionist program through the ``ASCH'' series (with Richard Lyons and Ronald Solomon) to produce a coherent, corrected, and streamlined proof.

\section{How It Was Solved and the Intuitions/Tools Needed}

\subsection{The Strategy of the Classification}

The classification strategy is based on the concept of the \textbf{centralizer of an involution}. The idea is:
\begin{enumerate}
    \item Every finite simple group contains involutions (elements of order 2) by Cauchy's theorem and the Feit--Thompson theorem (odd order groups are solvable).
    \item Study the centralizer of an involution, $C_G(z)$, which has a structure that can be analyzed.
    \item Show that the possible centralizer structures are limited and each corresponds to a known simple group.
\end{enumerate}

\subsection{Aschbacher's Key Contributions}

\begin{itemize}
    \item \textbf{Quasithin Groups:} A group $G$ is \textbf{quasithin} if its 2-rank is at most 2. These groups had been partially analyzed but remained a significant gap. Aschbacher and Smith produced a complete classification in two volumes (2004), resolving one of the last remaining pieces of the classification.
    \item \textbf{Groups of Characteristic 2-Type:} A group has \textbf{characteristic 2-type} if all 2-local subgroups have a strongly embedded 2-local structure. Aschbacher developed the signalizer functor method and the theory of pushing-up to analyze these groups.
    \item \textbf{The Uniqueness Case:} Aschbacher's theorem on uniqueness cases eliminated the possibility of unknown simple groups with certain local structural properties.
\end{itemize}

\subsection{Tools and Techniques}

\begin{itemize}
    \item \textbf{Signalizer Functors:} A method for finding normal subgroups in groups with given centralizer structure.
    \item \textbf{The Bender Method:} Analyzing the structure of maximal subgroups and p-layers.
    \item \textbf{Amalgam Method:} Studying the structure of groups acting on graphs to understand local structure.
    \item \textbf{The Pushing-Up Technique:} A method to control the structure of p-local subgroups.
    \item \textbf{Aschbacher's Block Theory:} A classification of subgroups of classical groups.
\end{itemize}

\section{Mathematical Theory}

\subsection{Finite Simple Groups}

\begin{definition}[Simple Group]
A group $G$ is \textbf{simple} if it has no proper nontrivial normal subgroups. That is, the only normal subgroups of $G$ are $\{e\}$ and $G$ itself.
\end{definition}

\begin{definition}[Classification of Finite Simple Groups (CFSG)]
Every finite simple group is isomorphic to:
\begin{enumerate}
    \item A cyclic group $\Z/p\Z$ of prime order;
    \item An alternating group $A_n$ for $n \geq 5$;
    \item A group of Lie type (Chevalley, Steinberg, Suzuki, Ree); or
    \item One of 26 sporadic simple groups.
\end{enumerate}
\end{definition}

\subsection{Aschbacher's Theorem on Maximal Subgroups}

\begin{theorem}[Aschbacher, 1984]
Let $G$ be a finite classical group (i.e., $\GL(n,q)$, $\SL(n,q)$, $\PSL(n,q)$, etc.) and let $H$ be a maximal subgroup of $G$ not containing the center of $G$. Then $H$ belongs to one of the following classes:
\begin{enumerate}
    \item $H$ stabilizes a proper subspace (reducible).
    \item $H$ stabilizes a direct sum decomposition (reducible).
    \item $H$ stabilizes a tensor product decomposition (tensor product).
    \item $H$ is defined over a subfield.
    \item $H$ normalizes an extension field subgroup.
    \item $H$ stabilizes a symmetric or alternating bilinear form.
    \item $H$ stabilizes a symmetric or alternating sesquilinear form.
    \item $H$ normalizes a classical subgroup.
    \item $H$ is almost simple and absolutely irreducible, not belonging to classes 1--8.
\end{enumerate}
\end{theorem}

This classification is fundamental for understanding the subgroup structure of classical groups and is widely used in group theory, number theory, and representation theory.

\subsection{The Aschbacher--Guralnick Theorem}

\begin{theorem}[Aschbacher--Guralnick, 1982]
Let $G$ be a finite simple group. Then $G$ can be generated by two elements. More precisely, for any nontrivial element $x \in G$, there exists $y \in G$ such that $G = \langle x, y \rangle$.
\end{theorem}

This result has important consequences for the structure of group extensions and the theory of group presentations.

\subsection{Fusion Systems}

\begin{definition}[Fusion System]
Let $G$ be a finite group and $S$ a Sylow $p$-subgroup of $G$. The \textbf{fusion system} $\mathcal{F}_S(G)$ is the category whose objects are subgroups of $S$ and whose morphisms are maps induced by conjugation by elements of $G$.
\end{definition}

Aschbacher's work on fusion systems has been influential in developing the abstract theory of saturated fusion systems, which generalize these categories to arbitrary groups and have applications in homotopy theory and algebraic topology.

\subsection{The Monster Group}

\begin{definition}[Monster Group]
The \textbf{Monster group} $\mathbb{M}$ is the largest sporadic simple group, with order:
\[
|\mathbb{M}| = 2^{46} \cdot 3^{20} \cdot 5^9 \cdot 7^6 \cdot 11^2 \cdot 13^3 \cdot 17 \cdot 19 \cdot 23 \cdot 29 \cdot 31 \cdot 41 \cdot 47 \cdot 59 \cdot 71
\]
\end{definition}

Aschbacher played a role in proving the uniqueness of the Monster and in establishing its existence through the study of group amalgams.

\section{Impact on Science and Technology}

\subsection{In Mathematics}

\begin{enumerate}
    \item \textbf{Group Theory:} The CFSG is the crowning achievement of 20th-century group theory. Aschbacher's contributions closed the most difficult cases.
    \item \textbf{Representation Theory:} Aschbacher's theorem on maximal subgroups is a cornerstone of modular representation theory.
    \item \textbf{Algebraic Combinatorics:} Group classification theorems have applications in graph theory, design theory, and coding theory.
    \item \textbf{Number Theory:} The classification of finite simple groups is used in class field theory and the theory of modular forms.
    \item \textbf{Algebraic Topology:} Fusion systems and their homotopy-theoretic analogues (p-local groups) draw on Aschbacher's work.
\end{enumerate}

\subsection{In Other Fields}

\begin{enumerate}
    \item \textbf{Cryptography:} Finite group theory is used in cryptographic systems. The structure theory of groups of Lie type informs elliptic curve cryptography and isogeny-based cryptography.
    \item \textbf{Physics:} Sporadic groups (particularly the Monster) appear in the study of vertex operator algebras and conformal field theory.
    \item \textbf{Computational Group Theory:} The classification provides the foundation for algorithms in computational group theory (GAP, Magma).
\end{enumerate}

\section{Five Frequently Asked Questions}

\subsection*{Q1: What is the Classification of Finite Simple Groups?}
The CFSG is a theorem stating that every finite simple group belongs to one of four types: cyclic of prime order, alternating, Lie type, or one of the 26 sporadic groups. It is considered one of the largest and most complex theorems ever proved.

\subsection*{Q2: What did Aschbacher specifically contribute?}
Aschbacher solved the quasithin group case (the last major gap), developed the theory of signalizer functors and pushing-up, classified maximal subgroups of classical groups, and was the leading architect of the revisionist program.

\subsection*{Q3: Why is the classification important?}
It provides a complete description of the building blocks of all finite groups. Every question about finite groups can, in principle, be reduced to questions about the simple groups, whose structure is now fully known.

\subsection*{Q4: What are sporadic groups?}
The 26 sporadic simple groups are those that do not fit into any infinite family. Examples include the Mathieu groups, the Janko groups, the Conway groups, the Fischer groups, and the Monster. The largest, the Monster, has about $8 \times 10^{53}$ elements.

\subsection*{Q5: What should an undergraduate study to understand Aschbacher's work?}
Group theory (through a standard graduate sequence), representation theory of finite groups, and Lie theory. Recommended: \textit{Finite Group Theory} (Aschbacher), \textit{The Classification of the Finite Simple Groups} (Gorenstein--Lyons--Solomon), and \textit{Lectures on Chevalley Groups} (Steinberg).

\section{Related Chapters}

See also Chapter~26 (Thompson) on finite simple group classification.

\section*{Exercises for the Reader}
\addcontentsline{toc}{section}{Exercises}

\begin{enumerate}
    \item [B] Prove that $A_5$ is simple by analyzing its conjugacy classes.
    \item [I] Show that the order of the Monster group is divisible only by the primes listed above.
    \item [B] Classify all groups of order less than 60.
    \item [A] Prove Aschbacher's theorem for $\GL(2,q)$ explicitly (the maximal subgroups).
    \item [I] Compute the 2-rank of $\PSL(2,q)$ for various $q$.
    \item [I] Show that the centralizer of an involution in $A_n$ is isomorphic to a wreath product $S_2 \wr S_{n-2}$ for $n \geq 5$.
    \item [A] Prove that every finite simple group is generated by two elements, using the classification.
\end{enumerate}

\section*{Timeline}
\addcontentsline{toc}{section}{Timeline}

\begin{tabularx}{\linewidth}{|l|X|}
\hline
\textbf{Year} & \textbf{Event} \\ \hline
1944 & Born in Little Rock, Arkansas \\ \hline
1966 & B.A.\ from Princeton University \\ \hline
1969 & Ph.D.\ from the University of Chicago (advisor: Jonathan Alperin) \\ \hline
1971 & Completes proof of the uniqueness case for groups of characteristic 2-type \\ \hline
1973 & Moves to Caltech as Professor (remains ever since) \\ \hline
1976 | Solves the quasithin problem (first of three main solutions) \\ \hline
1980 | Cole Prize of the AMS \\ \hline
1984 | Publishes Aschbacher's theorem on maximal subgroups \\ \hline
1990 | Rolf Schock Prize \\ \hline
2004 | Publication of \textit{The Classification of Quasithin Groups} (with Smith) \\ \hline
2011 | Publication of \textit{Fusion Systems in Algebra and Topology} \\ \hline
2012 | Wolf Prize in Mathematics (with Lyons, Smith, Solomon) \\ \hline
\end{tabularx}

\section*{Glossary}
\addcontentsline{toc}{section}{Glossary}

\begin{description}
    \item[Amalgam] A group formed by identifying subgroups of two groups along a common subgroup.
    \item[Characteristic $p$-type] A group in which all $p$-local subgroups are $p$-constrained.
    \item[Fusion system] A category encoding the conjugacy structure among $p$-subgroups of a group.
    \item[Involution] An element of order 2 in a group.
    \item[$p$-local subgroup] The normalizer of a nontrivial $p$-subgroup.
    \item[Quasithin group] A finite group whose 2-rank is at most 2.
    \item[Signalizer functor] A consistent assignment of subgroups to elementary abelian subgroups, used to find normal subgroups.
    \item[Simple group] A group with no proper nontrivial normal subgroups.
    \item[Sporadic group] A finite simple group not belonging to any infinite family.
    \item[Sylow $p$-subgroup] A maximal $p$-subgroup of a finite group (exists by Sylow's theorems).
\end{description}

\section{Citations}

\subsection*{Primary Sources}
\begin{enumerate}
    \item M.\ Aschbacher, \textit{Finite Group Theory}, Cambridge Univ.\ Press, 1986.
    \item M.\ Aschbacher, \textit{On the maximal subgroups of the finite classical groups}, Invent.\ Math.\ 76 (1984), 469--514.
    \item M.\ Aschbacher, \textit{A condition for the existence of a strongly embedded subgroup}, Proc.\ AMS 38 (1973), 509--511.
    \item M.\ Aschbacher and R.\ Guralnick, \textit{Some applications of the first cohomology group}, J.\ Algebra 90 (1984), 446--460.
    \item M.\ Aschbacher and S.\ D.\ Smith, \textit{The Classification of Quasithin Groups}, AMS Math.\ Surveys and Monographs, 2004.
    \item M.\ Aschbacher, R.\ Lyons, and S.\ D.\ Smith, \textit{The Classification of Finite Simple Groups: Groups of Characteristic 2-Type}, AMS, 2011.
\end{enumerate}

\subsection*{Historical References}
\begin{enumerate}
    \item D.\ Gorenstein, \textit{Finite Simple Groups}, Plenum Press, 1982.
    \item D.\ Gorenstein, R.\ Lyons, and R.\ Solomon, \textit{The Classification of the Finite Simple Groups}, AMS Math.\ Surveys and Monographs (multiple volumes, 1994--present).
    \item R.\ Solomon, \textit{A brief history of the classification of the finite simple groups}, Bull.\ AMS 38 (2001), 315--352.
\end{enumerate}

\subsection*{Expository Works}
\begin{enumerate}
    \item M.\ Aschbacher, \textit{The status of the classification of the finite simple groups}, Notices AMS 51 (2004), 736--740.
    \item M.\ Aschbacher, \textit{The classification of the finite simple groups}, in \textit{The Mathematical Heritage of the Twentieth Century}, 1991.
    \item J.\ H.\ Conway, R.\ T.\ Curtis, et al., \textit{Atlas of Finite Groups}, Clarendon Press, 1985.
\end{enumerate}

